% !TeX root = ../main.tex
%------------------------------------------------------------------------------
%  Three-Dimensional Ground Modelling
%------------------------------------------------------------------------------
\chapter{Three-Dimensional Ground Modelling}
\label{ch:ground-modelling}

Every subsequent stage of design draws on the ground model.
Empirical classification requires domains, analytical solutions require a
stratigraphy and a stress state, numerical analysis requires geometry and
material assignment, and construction control requires an expectation against
which the exposed face can be compared. The model is therefore the principal
engineering deliverable of an investigation rather than a by-product of
it~\citep{knill2003ground}.

The way the model is built has changed. Hand-drawn sections interpolated
between boreholes have given way to implicit surfaces fitted to all available
data, updated as each new borehole, geophysical line or mapped face arrives,
and increasingly conditioned by statistical and machine-learning methods. The
gain is that the model can be rebuilt cheaply
enough to be kept current, and that its uncertainty can be quantified rather
than implied by the confidence of the draughtsman.

This chapter sets out how such a model is built and kept current. It first
distinguishes the models a project carries and the data that support them. It
then treats the construction of the geometry, the division into domains and
the assignment of properties, and the description of uncertainty. It closes
with the updating of the model through investigation and construction and its
handover as a project record.

\section{From data to ground model}
\label{sec:data-to-ground-model}

\subsection{Model types and their purposes}
\label{sec:model-types}

Four objects are commonly conflated: the three distinguished in
Section~\ref{sec:ground-model} and a hydrogeological model. The geological model interprets units,
contacts and structures. The hydrogeological model adds aquifers, aquitards,
heads and boundaries, and need not share the geological model's unit
boundaries: a single lithological unit may contain several flow domains,
and a permeable fault may connect units that are otherwise separated. The
ground model attaches material behaviour and design parameters.
The design model is the idealisation actually meshed and solved.

Each has a different tolerance. A geological model that misplaces a contact by
five metres may be adequate for route selection and useless for
setting a segment ring's design load. Building all four as a single object,
which modern software encourages, hides these differing tolerances and
transfers the finest requirement to every user.

\subsection{Data sources and their support volumes}
\label{sec:data-support}

Data differ in what volume of ground each observation represents. A borehole
samples a cylinder tens of millimetres across; a laboratory specimen a volume
smaller again. A seismic refraction line averages properties over tens of
metres. A mapped exposure or tunnel face gives dense two-dimensional structural
information over a plane but says nothing directly about the volume behind it.
Machine data from a tunnel boring machine integrates the response of the full
face over one advance of the machine, typically one to two metres
(Chapter~\ref{ch:tbm}).

Combining these without regard to support volume produces a model more precise
than the data justify. A point strength measurement and a geophysical
velocity are not competing estimates of the same quantity: the first is a
sample of a specific piece of rock, the second an average that includes the
fractures the first avoided. The usual reconciliation is to use the dense but
indirect data to define geometry and the sparse but direct data to assign
properties within it.
Where a correlation can be established and tested against withheld data, the
indirect record can carry properties as well as geometry
(Section~\ref{sec:ml-ground-modelling}).

Sampling bias compounds this. Boreholes are drilled where access permits and
where trouble is expected, so a borehole population is rarely a random sample
of the alignment. Corrections exist for orientation bias in discontinuity
sampling (Chapter~\ref{ch:structural-analysis}, Section~\ref{sec:statistical-treatment-of-orientation-data}); there is no comparable
correction for having sampled only the ground that was easy to reach.

\subsection{Model resolution and decision scale}
\label{sec:model-resolution}

Resolution should follow the decision. Route selection needs formation-scale
geometry over kilometres. Reference design needs unit boundaries to a few
metres and parameters by domain. Detailed design needs the geometry of
individual features that govern a specific structure. Construction control
needs metre-scale ground classification updated at the face. These are
different requirements, and satisfying the finest everywhere is expensive and
misleading. Figure~\ref{fig:support-scale} sets them against the support of
the data.

\begin{figure}[htbp]
\centering
\datafigure[\textwidth]{ch06/support-scale}
\caption{Characteristic length represented by each data source against the resolution each decision requires; indicative ranges from Sections~\ref{sec:data-support} and~\ref{sec:model-resolution}. No source matches every decision, and point measurements are one to two orders of magnitude finer than any decision they inform.}
\label{fig:support-scale}
\end{figure}

The risk of misleading is the more serious. A ground surface rendered
to millimetre tolerance between boreholes two hundred metres apart reflects
the interpolator rather than the ground, and it will be read as
survey by anyone who did not build it. Resolution should be declared with the
model, and where it varies along the alignment, as it always does, that
variation should be visible in the deliverable.

\section{Building the geometry}
\label{sec:building-geometry}

\subsection{Explicit and implicit modelling}
\label{sec:implicit-modelling}

Explicit modelling constructs surfaces from digitised sections and hand-drawn
contacts: the geologist decides where the boundary runs and the software
records it. Implicit modelling instead fits a scalar field to the observations
(contact points, above-and-below constraints, structural orientation
measurements) and extracts the boundary as an isosurface of that field.

The practical difference is reproducibility and update cost. An implicit model
is a function of its inputs, so adding a borehole regenerates the surface in
minutes and the change is attributable to the new datum
(Figure~\ref{fig:implicit-scalar-field}). An explicit model
embodies judgement that is not recorded anywhere except in the geologist's
head, and updating it means redrawing.

\begin{figure}[htbp]
\centering
\datafigure[\textwidth]{ch06/implicit-scalar-field}
\caption{Implicit modelling of a contact (schematic section). A scalar field (grey isolines) is fitted to contact points, above-and-below constraints along each borehole and one dip measurement, and the contact is its zero isoline. A fifth borehole regenerates the surface; the change from the previous contact (dashed) is largest near the new datum.}
\label{fig:implicit-scalar-field}
\end{figure}

The cost is a loss of control. An interpolator has no geological knowledge, and left
unconstrained it will produce surfaces that are smooth, plausible-looking and
geologically impossible: units that pinch and reappear without cause,
contacts that cross, folds with no consistent axis. Implicit modelling
therefore records a geological interpretation and does not replace it: the
geologist's contribution moves from drawing the surface to
supplying the constraints that keep the surface geologically admissible.

\subsection{Structural and stratigraphic control}
\label{sec:structural-control-modelling}

The constraints that matter are the geological rules. Stratigraphic order must
be imposed so that units appear in sequence and an erosional surface truncates
what lies beneath it. Fault displacement must be applied so that the modelled
offset matches the observed sense and magnitude rather than being smeared into
a gradual flexure. Fold geometry is best handled by supplying orientation data
(measured dips, not just contact points) so the interpolator reproduces
limb dips instead of averaging across the hinge. Intrusive relationships must
be declared, since an intrusion cross-cuts rather than conforms.

Where these rules are supplied, the model reproduces the geologist's
understanding and generalises it consistently between control points. Where
they are not, the model reproduces the mathematics of the interpolator.
Reviewing a ground model therefore means reviewing its constraints, and a
model delivered without them cannot be reviewed.

\subsection{Faults, domains and discontinuity networks}
\label{sec:dfn}

Faults enter the model as discrete volumes, bounded by surfaces and carrying
an assigned damage-zone thickness, because their engineering effect (weakness, deformability, permeability,
abrupt class change) is distributed across the damage zone rather than
concentrated on a plane. Modelling a fault as a zero-thickness surface makes
it invisible to any subsequent volume-based calculation.

Below fault scale, discontinuities are too numerous to model individually and
too influential to omit. A discrete fracture network generates a synthetic
population from the statistics established in
Chapter~\ref{ch:structural-analysis}: orientation distributions by set,
spacing or intensity, persistence, and termination rules. The result is one
realisation of a population with the observed statistics, and not a map of
the actual fractures.

That distinction governs how such networks may be used. They are appropriate
for questions about the distribution of outcomes: the range of block sizes,
the probability of an adverse wedge on a given face, the connectivity that
controls flow. They are not appropriate for predicting what will be found at a
specific chainage. Conditioning a network to observed intersections, where
mapping or borehole data allow, narrows the gap but does not close it.

\section{Populating the model}
\label{sec:populating-model}

\subsection{Assigning parameters and domains}
\label{sec:domain-assignment}

Before properties are distributed, the model is divided into geotechnical domains: volumes within which
behaviour is sufficiently uniform that one parameter set applies. Domain
boundaries follow whichever control dominates (lithology, weathering grade,
structural intensity, alteration), and they need not coincide with
geological unit boundaries.

The parameter derivation chain of
Chapter~\ref{ch:rock-mass-parameters} should be carried out inside the model
rather than in a separate spreadsheet, for traceability. When
$\sci$, $\mi$, $\GSI$ and $\Dfac$ are attributes of a domain, the derived
$\mb$, $s$, $a$ and $\Erm$ update automatically when the domain is revised,
and the value used in an analysis can be traced to the observations behind it.
When the derivation lives in a spreadsheet, the two diverge within weeks and
nobody can say which is current.

\subsection{Geostatistical estimation and simulation}
\label{sec:geostatistics}

Once geometry and domains exist, properties must be distributed within them. Geostatistics
begins by quantifying spatial
correlation~\citep{baecher2003reliability,phoon1999characterization} through
the variogram,
\[
\gamma(h) = \tfrac{1}{2}\,\mathrm{E}\!\left[\bigl(Z(x+h)-Z(x)\bigr)^{2}\right],
\]
which measures how rapidly the squared difference between two measurements
of a property $Z$ grows, on average, with their separation $h$. Three
quantities describe it (Figure~\ref{fig:variogram-model}): the sill it approaches, equal to the variance of the
field; the range at which it reaches the sill, beyond which two measurements
are effectively independent; and the nugget, the apparent variance as the
separation tends to zero, representing measurement error and variability at
scales below the sample spacing.

\begin{figure}[htbp]
\centering
\datafigure[0.62\textwidth]{ch06/variogram-model}
\caption{Spherical and exponential variogram models with a nugget of 0.2 and a range of 100~m, normalised by the field variance $\sigma^2$ so that the sill is unity. The exponential model approaches the sill asymptotically and reaches about 95 per cent of it at the range.}
\label{fig:variogram-model}
\end{figure}

Kriging then estimates values at unsampled locations as a weighted average of
the observations, with weights derived from the variogram. It is optimal in a
least-squares sense, and it is smooth: it honours the data at the observation
points, but away from them the estimate reverts toward the mean, and the
variance of the estimated field is lower than the variance of the real one
(Figure~\ref{fig:kriging-simulation}).

Because of that smoothing, estimation is often unsuitable.
Underground design is frequently governed by extremes (the weakest zone,
the most permeable path), and a smoothed field has removed them.
Conditional simulation instead generates multiple realisations, each honouring
the data at observation points and reproducing the observed variogram in
between. No single realisation is the best estimate, but the ensemble
describes the variability the design must tolerate. Where a design depends on
variability rather than on the mean, simulation is the appropriate method and
kriging is not.

\begin{figure}[htbp]
\centering
\datafigure[\textwidth]{ch06/kriging-simulation}
\caption{Kriging and conditional simulation of GSI along a 1~km alignment from six boreholes (synthetic field: spherical variogram, range 150~m, standard deviation 8). The kriged estimate reverts to the mean between holes and has about half the variability of the true field; each realisation honours the boreholes and reproduces that variability. The band is the kriged estimate $\pm 2$ kriging standard deviations.}
\label{fig:kriging-simulation}
\end{figure}

Both estimation and simulation assume that the statistics are stationary over
the volume to which they are applied, that is, that one variogram describes
it. Ground rarely satisfies this. A variogram fitted across a fault or a weathering
front reports a range and a sill that belong to no part of the ground, and
the realisations built on it are wrong everywhere. Geostatistics is therefore
applied inside the domains of Section~\ref{sec:domain-assignment} and not
across them, which is the same requirement that geological constraint
imposes on the geometric interpolator.

\subsection{Machine learning for property and lithology prediction}
\label{sec:ml-ground-modelling}

Supervised learning is useful where a property is expensive to
measure and correlated with something cheap and abundant. Downhole geophysical
logs, core photography and drilling
parameters~\citep{schunnesson1996rqd} are all recorded continuously
or nearly so; lithology, weathering grade and strength are determined at
intervals. Training a classifier on the logged intervals and applying it to
the unlogged remainder converts sparse determinations into a continuous
profile.

The technique works under three conditions. Training data must span the
classes to be predicted, and the classes are almost always imbalanced: the
fault gouge that matters most is the rarest label in the set. Validation must
be against withheld boreholes rather than withheld samples, because samples
from the same borehole are not independent and a model validated against them
will report an accuracy it does not have. Finally, the predicted quantity must
be correlated with the inputs and not only co-located with them.

Automated interpretation of core imagery deserves separate mention because it
addresses a bottleneck in logging. Fracture counting and $\RQD$ determination are
tedious, subjective and slow, and consistency between loggers is poor. Image
analysis is repeatable, which is not the same as being correct, but
repeatability is valuable when the output feeds a
classification system calibrated on other people's logging.

\section{Uncertainty}
\label{sec:model-uncertainty}

\subsection{Sources of model uncertainty}
\label{sec:model-uncertainty-sources}

Model uncertainty comes from four places: the density and spatial distribution
of the data; ambiguity in the geological interpretation; the choice of
interpolation method and its parameters; and the intrinsic complexity of the
geology being modelled. These are not equally reducible. More drilling
addresses the first and, sometimes, the second. The third is reducible without
any new data: an implicit model is cheap to rebuild, so the
interpolator and its parameters can be varied deliberately and the spread of
the results reported. Nothing addresses the fourth.

None of this is visible in the output. A model built from
three boreholes and one built from three hundred render identically, in the
same colours, with the same crisp surfaces. Every subsequent user sees the
same confident appearance and none of the evidence behind it.

\subsection{Representing and communicating uncertainty}
\label{sec:uncertainty-representation}

Several devices restore the missing information. Geometry can be simulated
as well as estimated, by perturbing the implicit scalar field within the
tolerance its constraints allow or by categorical simulation of the unit
indicator, and the realisations referred to here are the output of that
step. Multiple realisations,
displayed as an ensemble rather than as a single surface, show directly how
much the boundary could move. Probability-of-lithology volumes report, for
each location, the proportion of realisations assigning it to each unit,
which makes the transition zones visible as zones rather than as lines
(Figure~\ref{fig:contact-probability}).
Distance-to-nearest-datum shading is cruder but immediately legible, and it
shows what a reviewer most often wants to know.

\begin{figure}[htbp]
\centering
\datafigure[\textwidth]{ch06/contact-probability}
\caption{Probability of rock below a contact, from 200 conditional simulations of its elevation through five boreholes (synthetic; vertical scale exaggerated). Heavy line: kriged contact; grey: six realisations; dashed: 5 and 95 per cent probability. Lower panel: standard deviation of the estimated contact, classed as high ($<1$~m), moderate or low ($>3$~m) confidence.}
\label{fig:contact-probability}
\end{figure}

Simplest and most useful is a confidence classification along the alignment:
high, moderate or low, with the criteria stated. It carries into the
geotechnical baseline report~\citep{essex2007gbr}, into investigation
planning, and into the
construction-stage expectation of how closely the ground should match the
prediction. Where the model says confidence is low and the ground is
consistently as predicted, that is worth knowing; where confidence is high and
the ground differs, something is wrong with the model rather than with the
ground.

\subsection{Value of information}
\label{sec:value-of-information}

Investigation budgets are finite, and uniform borehole spacing is almost never
the best use of them. The value of an additional observation is the extent to
which it changes a decision~\citep{baecher2003reliability}, and that is not
distributed uniformly along an alignment.

The reasoning is straightforward even when the formal calculation is not; the
formal version, with prior, likelihood and posterior stated, is given in
Section~\ref{sec:bayesian-updating}. An
additional borehole is worth drilling where uncertainty is high \emph{and} the
consequence of being wrong is high \emph{and} the result could plausibly
change what is done. Uncertainty alone is insufficient: there is no value in
resolving a boundary in a reach where the design is the same either way.
Consequence alone is insufficient: there is no value in re-confirming ground
that is already well characterised.

Used this way the ground model becomes an investigation planning tool as much
as a design input, and the investigation becomes a sequence of decisions
rather than a fixed programme executed to completion.

\section{Keeping the model current}
\label{sec:model-updating}

\subsection{Updating during investigation and construction}
\label{sec:model-update-cycle}

The model should be regenerated whenever new data arrive: boreholes and
geophysics during investigation, then probe holes, mapped faces, machine data
and instrument readings during construction. The construction phase generates
far more ground information than the investigation did, and it arrives
progressively, at the face, in time to influence what is done next.

Sequential updating in a Bayesian sense, treating the existing model as a
prior and the new observation as evidence, is the appropriate formalism,
and the informal version is universally practised even where the mathematics
is not. What distinguishes a disciplined process is that updates are recorded,
attributed and dated, so that a later reader can see how the interpretation
evolved and on what evidence.

\subsection{Reconciling prediction with observation}
\label{sec:prediction-vs-observation}

Systematic comparison of predicted against encountered conditions, chainage by
chainage, serves three purposes. It is the evidence base for baseline
comparison and therefore for claims. It calibrates the model, showing whether
the interpretation is biased in a consistent direction. It also supports
revising the ground class expectation ahead of the face, which is
how observational design
operates~\citep{peck1969observational,nicholson1999observational}.

The comparison is only meaningful if the prediction was recorded before the
ground was seen. This requirement is routinely violated: a ground class
assigned at the face and entered into the same record as the prediction
destroys the evidence. Prediction and observation must be separately recorded
and separately dated.

\subsection{The model as a project record}
\label{sec:model-as-record}

An operator inheriting a tunnel needs to know what ground it passes through,
what was encountered during construction, and where the interpretation was
uncertain. That information exists at handover and is routinely lost, because
it lives in a proprietary model file that no one maintains and in a set of
construction records that are archived rather than integrated.

Preserving it requires three things: version control with the provenance of
every surface and parameter recorded~\citep{iso19650}; export in an open format that will
outlive the software (Chapter~\ref{ch:digital-twin}); and an explicit handover
of the model as a deliverable in its own right, with a named recipient. The
asset will be inspected, assessed and possibly modified for a century. The
ground model is the document that explains what the asset was built in, and
without it neither the construction records nor the design reports can be
read against the ground they describe.

\chapterreferences
